% !TeX root = ../main.tex
\halamanAwal{KATA PENGANTAR}
\begingroup
\setstretch{1.11}
\setlength{\parskip}{5pt}
\setlength{\emergencystretch}{2em}

Puji syukur penulis panjatkan ke hadirat Tuhan Yang Maha Esa atas berkah dan rahmat-Nya sehingga usulan penelitian (proposal tugas akhir) yang berjudul “\judulIndonesia” dapat terselesaikan dengan baik. Naskah proposal tugas akhir ini disusun sebagai salah satu syarat untuk menempuh Seminar Proposal pada Program Studi Sains Data, Fakultas Sains, Institut Teknologi Sumatera. Penyusunan naskah proposal ini tidak lepas dari bimbingan, arahan, dan dukungan berharga dari berbagai pihak. Dalam kesempatan ini, dengan penuh rasa hormat dan kerendahan hati, penulis menyampaikan ucapan terima kasih yang sebesar-besarnya kepada:

\begin{enumerate}
  \setlength{\itemsep}{3pt}
  \setlength{\parsep}{0pt}
  \item \pembimbingSatu{} selaku Dosen Pembimbing I dan \pembimbingDua{} selaku Dosen Pembimbing II atas waktu, bimbingan, dan arahan yang diberikan selama penyusunan proposal ini.
  \item Bapak/Ibu Dosen Penguji yang telah memberikan masukan, evaluasi, dan saran perbaikan yang konstruktif demi kesempurnaan penelitian ini.
  \item \koordinatorProdi{} selaku Koordinator Program Studi \namaProgramStudi{} atas dukungan dan arahan akademik selama penyusunan proposal ini.
  \item Seluruh Bapak/Ibu Dosen dan Tenaga Kependidikan di lingkungan Program Studi Sains Data, Fakultas Sains, Institut Teknologi Sumatera atas ilmu, bimbingan, dan fasilitas akademik yang telah diberikan.
  \item Kedua orang tua dan keluarga tercinta atas doa, motivasi, kasih sayang, dan pengorbanan yang tiada henti.
  \item Rekan-rekan mahasiswa Program Studi Sains Data ITERA angkatan 2023 dan seluruh pihak yang telah memberikan bantuan teknis, diskusi, serta semangat dalam penyelesaian proposal ini.
\end{enumerate}

Penulis menyadari bahwa naskah proposal ini masih memiliki keterbatasan. Oleh karena itu, penulis sangat mengharapkan kritik dan saran yang membangun dari para penguji dan pembaca guna penyempurnaan pelaksanaan penelitian ke depan. Semoga penelitian ini dapat memberikan kontribusi nyata bagi pengembangan ilmu pengetahuan di bidang rekayasa sistem basis data dan sains data.

\par\noindent
\begin{minipage}[t]{\linewidth}
  \raggedleft
  \namaLokasi, \tanggalDokumen\par
  \vspace{0.9cm}
  \namaPenulis\par
\end{minipage}\par
\clearpage
\endgroup
